\section{Reference Pipeline Implementation}

\subsection{Foundation, configuration, and operational controls}

The early scaffold was converted into an installable Python package with a
neutral Docker entry point and one authoritative dependency file. Its command
line separates preprocessing, modelling, calibration, scoring, adaptation,
evaluation, and benchmarking into explicit stages. Commands return structured
output, use stable error codes, support dry runs that validate configuration without executing the pipeline, and write UTC
logs.

Each device has a strict JSON configuration. It records the device identity and
MAC address, input and output locations, protocol order, window length,
chronological partitions, estimator, likelihood backend, calibration rule,
adaptation mode, and random seeds. Misspelled fields, inconsistent category
counts, invalid quantiles, duplicate capture assignments, reversed partitions,
and impossible adaptation buffer sizes are rejected before a run begins. The
shell experiment runner selects reviewed configuration files instead of
requiring detector edits for each run.

Pydantic schemas validate models, thresholds, events, and manifests when they
are loaded. Model and dataset fingerprints also prevent artefacts from different
runs from being combined merely because their filenames look compatible.

\subsection{Data engineering and capture validation}

Source PCAPs can first be inventoried, read completely for structural validation,
or fingerprinted without running model code. The fingerprint hashes both capture
bytes and the policies that give those bytes meaning, including category order,
windowing, and partition assignment. A changed file or experiment definition
therefore produces a different dataset identity.

The processor selects one device by Ethernet or ARP hardware address. It filters
raw frames before asking Scapy to decode them, an important improvement for large
mixed-device captures. An experiment on a 1.48~GB UNSW capture recorded a
reduction from 280.0 seconds to 12.5 seconds with byte-identical processed JSON.

Retained packets are classified once, timestamps are normalised to UTC, and
counts are placed in epoch-aligned half-open intervals $[t,t+10\text{ minutes})$.
Observed silence is stored as a valid zero vector. Missing capture coverage is a
different state and is excluded from model fitting and scoring. This distinction
prevents ``the device sent nothing'' from being confused with ``the recording
contains no evidence.''

The Chromecast and Samsung captures were filtered and converted into this common
window format. Their official attack annotations were stored separately using
the same capture ID and exact time bounds as the scored windows. Provenance
records document retained and excluded captures, incomplete days, file hashes,
attack intervals, and partition choices.

\subsection{Statistical modelling, calibration, and scoring}

Training loads only the declared fit captures, checks the device and category
names, and fits $\boldsymbol{\alpha}$ with Minka's fixed-point estimator. The
saved model identifies the source data, capture IDs, number of windows, and UTC
range. It also records the numerical backend, precision, fitted parameters,
whether fitting converged, the iteration count, the likelihood, and the runtime.

Calibration then scores a different benign partition. It stores the chosen
quantile, threshold, score minimum/median/maximum, and interquartile range. The threshold
also carries a deterministic model fingerprint. A model change invalidates the
old threshold rather than allowing an apparently compatible filename to pass.

Static scoring keeps the fitted model and threshold fixed. Every non-missing
development window produces an event containing its counts, score, decision,
severity, expected profile, category residuals, and model fingerprint. Residuals
identify which categories were high or low relative to expectation. They are a
local explanation, not an attack-family diagnosis.

The same pipeline also produces descriptive count statistics and separate daily
Dirichlet fits. Raw and per-packet-normalised score conventions are available as
explicit alternatives, but the reported detector uses raw scores throughout.
Because every event identifies its source data, model, and threshold, detector
variants can be compared window by window without guessing which artefacts were
used.

\subsection{Labelled evaluation}

Attack labels are stored separately from processed counts. Evaluation joins a
decision and label only when capture ID and exact time bounds agree. Duplicate
matches are rejected, and unmatched counts are reported. The evaluator produces
the confusion matrix, accuracy, precision, recall, F1, and false-positive rate. Labels are never
loaded during training, calibration, or adaptive updates.

The current \texttt{master} evaluator is deliberately limited to the development
partition. Final capture IDs exist in configuration, but the reference scoring
path has not yet been opened for a locked final run.

Recall is TPR, while FPR uses only the benign-label denominator:
\[
\mathrm{Precision}=\frac{TP}{TP+FP},\qquad
\mathrm{TPR}=\frac{TP}{TP+FN},\qquad
\mathrm{FPR}=\frac{FP}{FP+TN}.
\]
Accuracy is $(TP+TN)/(TP+TN+FP+FN)$ and F1 is the harmonic mean of precision and
recall. A zero precision/recall/FPR denominator produces JSON null; the current
F1 calculation also produces null when precision and recall sum to zero.
These conventions must be retained when presenting undefined metrics.

\subsection{Native LM integration}
\label{sec:lm-changes}

The professor-supplied code under \path{lm-algorithm/code/} is preserved as a
reference. The application calls a separate production copy under
\path{src/lm_idnet/algorithms/native/}. The mathematical approximation is
unchanged; the integration work makes it safe and efficient to use inside the
detector:

\begin{enumerate}
  \item Zero counts return the exact zero log-Gamma difference, so sparse and
  fully silent windows are valid inputs.
  \item Coefficient paths, file reads, and allocations are checked, and native
  failures are returned to Python instead of terminating the process.
  \item Native workspaces are released after successful, repeated, or failed
  calls.
  \item A matrix entry point validates dimensions and sums the likelihood of the
  independent rows used by the detector.
  \item The Python adapter converts fitted alpha parameters into the form
  expected by the native code, translates native errors into project exceptions,
  and protects the legacy global state.
  \item Setup that is common to every row is performed once for the complete
  matrix rather than repeated for every window.
  \item The container compiles the shared library reproducibly instead of using
  an unexplained pre-built object.
\end{enumerate}

The aggregate experiment that follows the reference paper checks convergence
with one pooled count vector, even though its estimator updates separate rows.
Since the Dirichlet--multinomial likelihood is nonlinear, the likelihood of the
pooled vector is not equal to the sum of the row likelihoods. The production
detector keeps the row-wise calculation and removes only repeated numerical
setup.

\subsection{The preserved legacy adaptation rules}

The following rules describe the preserved legacy policy. The conservative
default and its additional gates are described in Section~\ref{sec:conservative}.
Both use the same chronological scoring path. A configuration
selects \texttt{static}, \texttt{adaptive\_threshold}, or
\texttt{periodic\_refit}. In every mode the current window is scored before it can
affect future state:

\begin{center}
\fbox{score $\longrightarrow$ record decision $\longrightarrow$ update buffers
$\longrightarrow$ optionally promote future state}
\end{center}

Threshold-only adaptation keeps $\boldsymbol{\alpha}$ fixed and periodically
recalculates the lower quantile and IQR from a bounded buffer of recent scores.
It is intentionally naive: all recent non-missing scores enter the buffer,
including anomalous ones. A non-positive candidate IQR causes the update to be
rejected.

Periodic refitting also keeps a buffer of windows that appear safe. A window is
accepted only when its score under the original, unchanged detector satisfies

\[
s_0(\mathbf{x}_t)\geq
\tau_0+m\operatorname{IQR}_0.
\]

When enough accepted windows are available, the existing estimator tries to fit
a new alpha vector. Before fitting, one artificial single-count row is added for
each of the four protocol categories. This prevents a temporarily absent
protocol from receiving an invalid zero alpha. The detector accepts the new
model only if fitting converges, improves the likelihood of the accepted data,
and gives the recent scores a positive IQR. The alpha vector, expected protocol
profile, fingerprint, threshold, and IQR are updated together. If any check
fails, the old model remains active and only a threshold update is attempted.

The audit records when each update was attempted, the old and proposed values,
whether fitting converged and improved the likelihood, whether the update was
accepted, and why a rejection occurred. The saved original model and threshold
are never overwritten. The mechanism is therefore reproducible, but its legacy
safety checks do not prevent every form of model poisoning.

Static, threshold-only, and periodic-refit variants use the same implementation
and write to separate experiment directories. The initial experiments remain
available beside the later 288/240 development configuration, together with
reports of detection performance, buffer contamination, parameter movement, and
rejected updates.

\subsection{Validation and generated evidence}

The installed command is \texttt{lm-idnet}. Each subcommand performs one stage;
it never silently invokes a previous one. Every command also supports \texttt{--dry-run}; it does not execute the stage
or alter detector state, although CLI logging may create a log file. Expected failures are converted into stable
typed errors and non-zero exit codes, while UTC logs provide an operational
record.

At this report boundary the complete suite passes \textbf{560 tests}, including
the optional Docker checks. The ordinary default selection excludes ten slow,
PCAP, or benchmark tests, and Docker checks require the image variable. The tests
cover configuration, timestamps, packet classification, aggregation, schemas,
failure handling, fitting, calibration, scoring, labelled evaluation, native LM
agreement, benchmarking, adaptation ordering, candidate promotion and rejection,
and output fingerprints.

\begin{longtable}{>{\raggedright\arraybackslash}p{0.22\textwidth}
>{\raggedright\arraybackslash}p{0.70\textwidth}}
\caption{Working command and evidence map.}\label{tab:command-map}\\
\toprule
\textbf{Command or location} & \textbf{Responsibility} \\
\midrule
\endfirsthead
\toprule
\textbf{Command or location} & \textbf{Responsibility} \\
\midrule
\endhead
\texttt{preprocess} & Inventory, validate, fingerprint, or convert configured
captures into device-specific windows. \\
\texttt{diagnose} & Report capture-level count statistics, dispersion, and daily
Dirichlet fits. \\
\texttt{train} & Fit and store the normal Dirichlet--multinomial model. \\
\texttt{calibrate} & Learn and store a model-bound benign threshold. \\
\texttt{score} & Run the frozen static detector on development traffic. \\
\texttt{adapt} & Run the configured static, threshold-only, or periodic-refit
system and write an adaptation audit. \\
\texttt{evaluate} & Align development events with labels and store detection
metrics. \\
\texttt{run} & Read bounded count-window JSON lines and flush shared detector events. \\
\texttt{rollback} & Restore the previous complete accepted model--threshold pair. \\
\texttt{benchmark} & Compare LM and SciPy using the same fit, calibration, and
scoring workload. \\
\texttt{forecast} & Run the current synthetic metric scaffold; this is not an
empirical forecasting result. \\
\path{artifacts/} & Models, thresholds, scored events, and experiment-specific
adaptive outputs. \\
\path{reports/} & Evaluation, adaptation, capture, fingerprint, and benchmark
records. \\
\path{logs/} & UTC application messages; run/rollback logs rotate, while research
command logs use their existing file logging. \\
\bottomrule
\end{longtable}
